\documentclass{article}
\usepackage[T1]{fontenc}
\usepackage{lmodern}
\usepackage{graphicx}
\usepackage{amsmath,amssymb}
\usepackage[a4paper,margin=2cm]{geometry}
\usepackage[absolute,overlay]{textpos}
\setlength{\TPHorizModule}{1mm}
\setlength{\TPVertModule}{1mm}
\usepackage[indonesian]{babel}
\usepackage{hyperref}
\setlength{\emergencystretch}{2em}
% unit: Halaman 91

\begin{document}

% === Halaman 91 (llm) ===

\begin{itemize}
    \item Contoh: Misalkan kriptanalis menemukan
    \begin{center}
        cipherteks $\text{C}$ dan plainteks berkorepsonden $\text{K}$ \\
        cipherteks $\text{E}$ dan plainteks berkorepsonden $\text{O}$.
    \end{center}

    \item Kriptanalis $m$ dan $n$ dari kekongruenan berikut:
    \begin{align}
        2 &\equiv 10m + b \pmod{26} && (\text{i}) \\
        4 &\equiv 14m + b \pmod{26} && (\text{ii})
    \end{align}

    \item Kurangkan (ii) dengan (i), menghasilkan
    \begin{align}
        2 &\equiv 4m \pmod{26} && (\text{iii}) \\
        \text{Solusi: } & m = 7
    \end{align}

    Substitusi $m = 7$ ke dalam (i),
    \begin{align}
        2 &\equiv 70 + b \pmod{26} && (\text{iv}) \\
        \text{Solusi: } & b = 10.
    \end{align}
\end{itemize}

\end{document}
